\section{Vérifier le régime étudié et les calculs de réponse}\label{sec:stationnarite}
Pour chaque graine $s$, bras et observable $X$, nous calculons
\[
 R_s(X)=\frac{\langle X\rangle_{3750<t\le4000}}
                  {\langle X\rangle_{3000<t\le3250}},\qquad
 t_X=\frac{\overline R(X)-1}{s_R/\sqrt{12}},
\]
où $s_R$ est l'écart-type des douze rapports. L'IC95 est
$\overline R\pm t_{0{,}975;11}s_R/\sqrt{12}$, avec
$t_{0{,}975;11}\simeq2{,}201$.
Ce diagnostic compare deux niveaux temporels éloignés ; il n'analyse
ni toute la structure de corrélation ni la stabilité de la distribution.

\begin{table}[htbp]\centering\small
\caption{Rapport des moyennes des derniers et des premiers 250 pas de la fenêtre terminale.}\label{tab:stationnarite}
\begin{tabular}{@{}lrrr@{}}
\toprule
Bras & $R(Y)$ (IC95) & $R(N)$ (IC95) & $R(\sum K)$ (IC95) \\
\midrule
Contrôle & $1{,}0029\pm0{,}0152$ & $1{,}0014\pm0{,}0132$ & $1{,}0043\pm0{,}0178$ \\
$A\times1{,}5$ & $0{,}9968\pm0{,}0139$ & $0{,}9979\pm0{,}0116$ & $0{,}9957\pm0{,}0167$ \\
$A\times1{,}5$, dotation compensée & $0{,}9987\pm0{,}0132$ & $0{,}9974\pm0{,}0122$ & $0{,}9998\pm0{,}0148$ \\
Entrantes $A\times1{,}5$ & $1{,}0088\pm0{,}0106$ & $1{,}0070\pm0{,}0090$ & $1{,}0092\pm0{,}0128$ \\
Entrantes $A\times0{,}75$ & $0{,}9906\pm0{,}0159$ & $0{,}9928\pm0{,}0139$ & $0{,}9905\pm0{,}0184$ \\
Entrantes $\gamma=0{,}6$ & $1{,}0097\pm0{,}0159$ & $1{,}0038\pm0{,}0133$ & $1{,}0100\pm0{,}0186$ \\
\bottomrule\end{tabular}\end{table}


Aucun de ces contrastes n'excède le seuil bilatéral usuel, mais les
intervalles autorisent des dérives de l'ordre du pour-cent. Les observations
temporelles d'un run ne sont pas utilisées comme 1000 répétitions indépendantes.
Le résultat ne vaut pas automatiquement pour une fenêtre antérieure,
une autre observable ou une configuration exploratoire non présentée.

\paragraph{Renouvellement et mémoire de l'état initial.}
Le renouvellement d'un décile fournit une autre échelle temporelle, mais
il faut préciser la grandeur de classement et distinguer la disparition
d'une entité de sa sortie du décile. Pour un ensemble $D(t_0)$ défini
à une date de référence, la rétention
$|D(t_0)\cap D(t_0+h)|/|D(t_0)|$ mesure la persistance dans le groupe,
tandis que la fraction de ses membres encore vivants mesure leur survie.
Ces deux courbes ne sont pas interchangeables ; le terme « premier
décile » doit préciser s'il désigne les valeurs basses ou les plus élevées.
Même une rétention faible ne démontre pas l'oubli de la configuration
initiale : des engagements et des corrélations peuvent persister au-delà
des individus. Ce contrôle doit être confronté aux autocorrélations des
observables et à des initialisations distinctes. Aucun temps d'oubli
quantifié n'est établi ici ; l'éloignement de la fenêtre finale et
l'absence de dérive détectée ne suffisent pas à le certifier.

Les bilans de population donnent exactement
$N(t+1)-N(t)=B(t)-D(t)$ si $B$ et $D$ désignent les entrées et sorties
du pas avec la même convention de comptage. Sur une longue fenêtre stable,
$\overline B\simeq\overline D$ et $\overline B\simeq\lambda$ ;
le rapport $\overline D/\overline N$ devient proche de
$\lambda/\overline N$. Ce n'est pas exactement la moyenne de
$D(t)/N(t)$, et les termes de bord ne sont pas nuls par décret.
La relation explique pourquoi une élasticité de mortalité peut être
presque l'opposée de celle de population sans constituer une découverte
causale indépendante.

L'élasticité de l'équation~\eqref{eq:eps} porte sur une intervention
finie. Sa moyenne et son IC sont calculés après formation des douze
log-rapports appariés, et non après moyenne de tous les points temporels
dans les deux bras. Les formules de décomposition sont sensibles au
choix de l'effectif : les entités présentes à la production peuvent
disparaître avant l'instantané de fin de pas. Cette différence est
conservée dans les données sources ; elle ne doit pas être effacée pour
forcer la fermeture numérique d'une identité mal définie.

\section{Ajustements de distributions : définitions et limites}\label{sec:ajustements}
\subsection{Masse, densité par classe et fonction de survie}
Pour une taille entière $S$, la masse $P(S=s)$ et la survie
$P(S\ge s)$ sont distinctes. L'histogramme représente ici la masse
d'une classe $[l,r[$ divisée par $r-l$, afin qu'une classe large ne
paraisse pas plus probable seulement parce qu'elle contient plus
de tailles. La courbe ajustée est sommée sur les mêmes tailles
entières avant cette division. Une pente de survie ne peut pas être
substituée directement à l'exposant de masse.

La vraisemblance descriptive de la loi~\eqref{eq:avalanche} est maximisée
par graine, sur toutes les tailles $s\ge1$. La normalisation numérique
utilise $1\le s\le20000$ ; la masse des mille dernières valeurs est
contrôlée inférieure à $10^{-10}$ à l'optimum. Les bornes numériques
de recherche sont $-2\le\alpha\le5$ et
$10^{-5}\le1/S_c\le2$ ; un optimum sur une borne signalerait une
identification insuffisante. Les 24 ajustements présentés sont intérieurs.
Le support d'ajustement n'est pas choisi pour maximiser la rectitude
apparente d'un sous-ensemble de points.

Les IC inter-graines et le bootstrap de runs entiers mesurent la
variabilité des estimateurs ou des courbes moyennes obtenus par ce
protocole. L'ajustement par produit de probabilités ne démontre pas
l'indépendance des avalanches d'un run. Nous n'en tirons donc pas
d'erreurs-types événementielles ni un test d'adéquation certifié.
Il faudrait notamment compléter par une comparaison de familles sur
le même support et un diagnostic d'adéquation absolue tenant compte
de la dépendance temporelle (Clauset, Shalizi et Newman (2009), références §\ref{sec:biblio}).

Le coefficient descriptif affiché est
\[
 R^2_{\log}=1-
 \frac{\sum_j[\log\widehat p_j-\log p_{j,\rm ajust}]^2}
      {\sum_j[\log\widehat p_j-\overline{\log\widehat p}]^2},
\]
sur les classes empiriques non vides. Ce n'est pas un pourcentage
de probabilité correctement prédit, et ce n'est pas le critère
optimisé par la vraisemblance. Il dépend du découpage des classes.
Les zéros ne sont pas remplacés par une petite valeur arbitraire
pour les rendre logarithmables.

\subsection{Durées : exponentielle discrète et choix de fenêtre}
Si $D\ge1$ suit une loi géométrique, avec probabilité d'arrêt $u$ à
chaque pas, alors $P(D\ge d)=(1-u)^{d-1}$ et
$\widehat u=1/\overline D$. Ainsi
$\widehat\beta=-\log(1-1/\overline D)$ dans
$P(D\ge d)=\exp[-\beta(d-1)]$.
Les régressions de la figure~\ref{fig:art_cycles} utilisent cette
forme normalisée, non une droite en log-log. Le $R^2$ y est calculé
sur les log-survies non nulles, aux durées 1 à 11. Les épisodes restent
potentiellement corrélés et les frontières de fenêtre tronquent certains
d'entre eux ; les IC sont donc fondés sur les cinq graines et non sur
un nombre fictif de récessions indépendantes.

\subsection{Le test de Vuong et ce qu'il compare réellement}
Dans les analyses existantes des valeurs nettes et intérêts, le seuil
$x_{\min}$ est choisi par minimisation d'une distance de
Kolmogorov--Smirnov pour une puissance continue. Au-dessus de ce même
seuil et sur les mêmes valeurs, on ajuste aussi une log-normale et
une exponentielle, conditionnées à $X\ge x_{\min}$.
La puissance comparée ici est \emph{sans coupure supérieure} : ce
n'est pas le modèle discret à coupure exponentielle des avalanches.
La troncature \emph{à gauche}, qui définit un échantillon de queue,
ne doit pas être confondue avec la coupure \emph{à droite}.

Pour les densités ajustées $f$ et $g$, on définit
$d_i=\log f(x_i)-\log g(x_i)$ puis
\[
 z=\frac{\sqrt n\,\overline d}{s_d},\qquad
 s_d^2=\frac1{n-1}\sum_i(d_i-\overline d)^2.
\]
Une valeur positive favorise $f$ en log-vraisemblance moyenne.
L'interprétation usuelle de $\pm1{,}96$ comme seuil à 5\,\%
nécessite des hypothèses de régularité, de séparation des modèles
et d'échantillonnage. Ici le seuil a été sélectionné, les entités
interagissent, et certaines log-normales ajustées sont très proches
d'une limite en puissance. Le score est donc rapporté comme diagnostic,
pas comme certification de ces hypothèses (Vuong (1989), références §\ref{sec:biblio}).

Le signalement $\sigma_{\rm LN}>10$ concerne 50 des 72 ajustements
d'intérêts et 46 des 72 ajustements de valeur nette. Il marque une
zone numérique, pas une frontière mathématique entre deux familles.
Deux groupes visibles dans un nuage de paramètres d'ajustement ne
signifient pas que les agents forment deux populations économiques :
chaque point de ce nuage représentait un ajustement, pas une entité.
Les résultats par ajustement et leurs identifiants sont disponibles dans
\chemin{data_revision/rev_vuong_points.csv}.

Comparer Pareto, puissance à coupure et log-normale demande de
fixer le même domaine, de décrire la sélection du seuil et d'évaluer
l'adéquation, pas seulement de compter les paramètres. Une coupure
qui divergerait avec la taille du système définirait une limite
intéressante, à vérifier sur plusieurs tailles ; elle ne rend pas
gratuit son ajustement à taille finie.


\section{Compatibilité numérique et granularité temporelle}
\subsection{Compatibilité des moteurs et changements de règles}
\label{sec:compatibilite}
Trois niveaux de contrôle doivent être distingués : conserver les
fonctionnalités du laboratoire de simulation, reproduire une trajectoire
dans un régime de référence, et comparer deux modèles dont les règles
diffèrent. Seul le deuxième porte directement sur la parité dynamique.

Le rapport M4.2, section « Parité M4B », documente des tests à
$A=1$, $\gamma=1/2$ et $k=2$ : identité des variables discrètes
et comparaison des flottants à tolérance relative $10^{-9}$.
La sonde de 1000 pas rapportée ne présente pas de divergence discrète,
avec un écart relatif maximal des séries de $3{,}1\,10^{-12}$.
Ce constat n'est pas une identité bit à bit ni une garantie à tout horizon.

Pour M4.4, dans le dossier de campagne
\chemin{m4_4_rebond_credit_soc/}, le fichier
\chemin{results/analysis/parity_deviations_8000.csv}
conserve 8000 écarts maximaux nuls sur les 26 colonnes comparées à
la réalisation M4.3 de référence, dans le régime homogène à cible
arithmétique. Le test \chemin{tests/test_parity_m4_3.py} définit ce
périmètre. Le test \chemin{tests/test_v1_equivalence.py} compare
séparément Live-v1 au mode de compatibilité
\texttt{richest\_lends}, avec une intervention hétérogène ; un contrôle
négatif vérifie que libérer le sens du prêt peut changer la trajectoire.
La présence de ce test définit une exigence vérifiable ; elle ne remplace
pas à elle seule un compte rendu d'exécution pour chaque version.

Les chemins de campagnes M4.4 sont conservés comme identifiants de
provenance. Le passage à une règle de marché différente doit être
distingué d'une extension d'instrumentation, sans renommer les archives
dont proviennent les données. La non-régression ne
permet donc pas de fusionner les échantillons M4B, Live et M4.4.

\subsection{Pourquoi le regroupement temporel n'est pas une symétrie}
\label{sec:composition_temps}
Un contre-exemple apparaît avant même d'introduire le marché. Sans bruit,
crédit ni naissances, posons $F(K)=(1-\delta)(K+AK^\gamma)$.
Deux pas successifs donnent
\[
F(F(K))=(1-\delta)^2(K+AK^\gamma)
 +(1-\delta)A\big[(1-\delta)(K+AK^\gamma)\big]^\gamma.
\]
Un pas regroupé, avec $A'=2A$ et
$1-\delta'=(1-\delta)^2$, donnerait au contraire
\[
G(K)=(1-\delta)^2(K+2AK^\gamma).
\]
En général $F\circ F\ne G$ : la seconde production se calcule sur
un stock déjà modifié. Par exemple, à $K=100$, $A=1$,
$\gamma=1/2$ et $\delta=0{,}01$, les valeurs sont environ
$118{,}142$ et $117{,}612$. Le marché ajoute une alternance de
rencontres, de contrats, d'intérêts et de faillites ; composer les lois
de Poisson ou les seuls chocs ne compose pas ces opérations couplées.

Les essais Live-v1 de la figure~\ref{fig:art_temps} mesurent la portée
collective de cette différence. Après conversion complète des paramètres,
le regroupement de deux pas donne, relativement à sa référence,
$N:1{,}067\pm0{,}018$, $K_{\rm tot}:1{,}158\pm0{,}030$ et
$Y:1{,}118\pm0{,}024$. En regroupant quatre pas, le rapport de
production vaut $1{,}346\pm0{,}026$. Le protocole part d'un état vide,
utilise trois graines et compare $1200<t_{\rm ancien}\le2000$ ;
les flux sont divisés par le facteur de regroupement, pas les stocks.
Les valeurs par graine et les autres variantes sont dans
\chemin{data_article/temps_graines.csv} et
\chemin{data_article/temps_points.csv}. Ces essais et le contre-exemple
établissent l'absence d'une symétrie exacte ; ils ne constituent pas
une étude de convergence vers une limite continue.

\figurine{art_temps}{1.0}{\textbf{Le changement de pas ne s'arrête pas à
trois substitutions.} Expériences Live-v1 depuis un état vide, trois
graines par configuration. Fenêtre commune $1200<t_{\rm ancien}\le2000$.
Les flux des pas regroupés sont divisés par $s$, pas les stocks. Les
trois premiers essais regroupent deux pas ; le quatrième en regroupe quatre.
Le dernier compare deux configurations à $\delta=0{,}002$ avant conversion,
et possède donc sa propre référence. Points : rapports par graine moyennés ;
barres : IC95 de Student, larges avec trois répétitions.}
